\documentclass[a4paper,12pt]{article}
\usepackage[T1]{fontenc}
\usepackage[utf8]{inputenc}
\usepackage[portuguese]{babel}
\usepackage{hyperref}

\title{Manual de Teoria Musical --- Análise com N-Gramas e NER}
\author{Jéssica Cunha \\ SPLN 2025/26}
\date{}

\begin{document}
\maketitle

\begin{abstract}
As três frases mais representativas do \textit{Manual de Teoria Musical} (LXPRO),
selecionadas com base num modelo de linguagem de bigramas com Laplace smoothing:

\begin{itemize}
  \item Fá maior tem um valor de -1 (um bemol).
  \item Começamos com Dó maior, que tem um valor de 0.
  \item Acorde de Sétima Dominante Quando um acorde é denominado simplesmente como um acorde de sétima, usualmente significa que é um acorde de sétima dominante.
\end{itemize}
\end{abstract}

\section{Fonte}
Texto extraído de:
\begin{itemize}
  \item LXPRO. \textit{Manual de Teoria Musical}. Documento PDF.
\end{itemize}

\section{Corpo do Manual}
Texto extraído do PDF com \texttt{pdftotext}. As entidades nomeadas estão anotadas
inline com \textbf{negrito} e etiqueta \textsc{[tipo]}, identificadas com spaCy
(\texttt{pt\_core\_news\_sm}).

\subsection*{Manual de Teoria Musical}
INTRODUÇÃO \textbf{A Música}\textsc{[MISC]} é uma manifestação artística e cultural que resulta da expressão de sentimentos através de sons e silêncios, sendo também uma linguagem lógica e científica, uma forma de comunicação universal. Os elementos principais da \textbf{Música}\textsc{[LOC]} são o ritmo, as alturas dos sons, a dinâmica e o timbre. Abaixo encontra-se uma tabela que mostra quais as matérias a ser abordadas em cada um dos níveis de formação musical -- tanto em termos de teoria, como de treino auditivo (em itálico).

A notação musical é constituída por todos os elementos empregados na música escrita, dividindo-se em 3 classes principais: sinais representativos da altura dos sons (pauta, notas e claves); da duração dos sons (figuras das notas) e da duração dos silêncios (figuras de pausa).

\textbf{Pauta Musical}\textsc{[MISC]} A música escreve-se na pauta musical, também denominada de \textbf{Pentagrama}\textsc{[LOC]} ou \textbf{Partitura}\textsc{[LOC]}, composta por 5 linhas horizontais paralelas e 4 espaços. Estas linhas e espaços chamamse linhas e espaços naturais, contando-se de baixo para cima. Para aumentar a extensão da pauta musical, recorre-se a pequenas linhas suplementares. Estas escrevem-se por cima ou por baixo da pauta, tomando respectivamente o nome de linhas suplementares superiores e inferiores. Estas linhas contam-se sempre a partir das cinco linhas principais que constituem a pauta. Ou seja, as linhas suplementares superiores contam-se de baixo para cima; as linhas suplementares inferiores contam-se de cima para baixo.

\textbf{Clave}\textsc{[LOC]} A clave é um símbolo que define a posição das notas na pauta e determina a altura dos sons que elas representam. A clave dá o seu nome à nota que lhe corresponder na mesma linha. Há vários tipos de clave: \textbf{Sol}\textsc{[PER]}, Fá e Dó. Diferentes claves são utilizadas para instrumentos diferentes. Por exemplo, as vozes graves usam geralmente a clave de Fá, enquanto que as mais agudas usam a clave de \textbf{Sol}\textsc{[PER]}. Alguns instrumentos usam mais do que uma clave, como é o caso do piano, utilizando duas pautas ligadas por uma chaveta. A mão direita geralmente toca na clave de \textbf{Sol}\textsc{[PER]} e a mão esquerda na clave de Fá. A altura de cada nota é representada pela sua posição na pauta em referência à nota definida pela clave utilizada, como mostrado abaixo. Por exemplo, a clave de \textbf{Sol}\textsc{[PER]} começa na segunda linha, pelo que a esta linha corresponde a nota \textbf{Sol}\textsc{[PER]}.

Ou seja, a \textbf{Clave de Sol}\textsc{[MISC]} coloca a nota \textbf{Sol}\textsc{[PER]} na segunda linha da pauta -- já que é nessa linha que a clave começa. Desta forma, poderemos por comparação descobrir as restantes notas. Assim, Dó encontra-se na primeira linha suplementar inferior, Ré no primeiro espaço suplementar inferior, \textbf{Mi na primeira linha}\textsc{[MISC]}, etc. De notar que, devido ao emprego das linhas suplementares, as notas podem subir e descer em altura muito para além do estabelecido pelas 5 linhas naturais. As notas graves tendem a ficar nas linhas mais abaixo; as agudas nas linhas mais elevadas. O mesmo princípio aplica-se às restantes claves. A clave atribui uma nota a uma determinada linha e, por comparação, encontram-se as outras notas -- ver figura acima para a exata localização das várias claves e respetivas notas. Na pauta, o primeiro símbolo a ser colocado é a clave, seguida da armação de clave (secção \textbf{Tonalidade}\textsc{[MISC]}) e fórmula de compasso (secção \textbf{Compassos}\textsc{[MISC]}). Só depois são colocadas as figuras musicais (secção \textbf{Figuras Musicais}\textsc{[MISC]}).

Outros \textbf{Símbolos}\textsc{[LOC]} \textbf{Vejamos}\textsc{[MISC]} alguns dos símbolos e termos utilizados em partituras na notação musical moderna: \textbf{Pauta}\textsc{[MISC]} ou \textbf{Pentagrama}\textsc{[LOC]} São 5 linhas e 4 espaços. A pauta musical serve para escrever as partituras (através de notas, pausas, claves, etc.). Linhas e espaços suplementares São linhas que existem acima ou abaixo da pauta, porque nem sempre as 5 linhas e 4 espaços são suficientes para receberem todas as notas da música. Linhas de compasso \textbf{Usadas}\textsc{[LOC]} para separar dois compassos.

Linha de compasso dupla Usada para separar duas secções da música.

Linha de compasso tracejada \textbf{Subdivide compassos}\textsc{[MISC]}.

\textbf{Barra final Marca}\textsc{[LOC]} o fim de uma composição.

Na música escrita, há duas características do som que são transmitidas através da pauta: a altura do som (que nos indica que notas estão a ser tocadas) e a sua duração. Como vimos, a altura dos sons (as notas) é definida pela posição de determinadas figuras na pauta. Essas figuras musicais são os símbolos que representam a duração das notas nessa mesma pauta, isto é, a duração do som a ser executado. As figuras são mostradas abaixo, por ordem decrescente de duração. Elas são: semibreve, mínima, semínima, colcheia, semicolcheia, fusa e semifusa. Cada figura tem metade da duração da anterior. Antigamente existia ainda a breve, com o dobro da duração da semibreve, a longa, com o dobro da duração da breve e a máxima, com o dobro da duração da longa, mas essas notas já não são usadas na notação atual. Existe também a quartifusa, figura raras vezes aplicada, tendo metade do valor da semifusa. A duração real (medida em segundos) de uma nota depende do andamento utilizado. Isso significa que a mesma nota pode ser executada com duração diferente em peças diferentes ou mesmo dentro da mesma música, caso haja uma mudança de andamento. Assim, a mesma figura dura mais tempo (em segundos) num andamento lento do que num andamento ligeiro. De realçar que não são apenas os sons que são representados nas partituras. As pausas são símbolos que representam o silêncio, isto é, o tempo em que o instrumento não produz qualquer som. As pausas subdividem-se, tal como as notas, em termos de duração. Cada pausa dura o mesmo tempo relativo que sua nota correspondente, ou seja, a pausa mais longa corresponde à duração de uma semibreve; a mais curta a uma semifusa.

\textbf{Ligadura}\textsc{[LOC]} de Prolongação / \textbf{Ponto de Aumentação}\textsc{[MISC]} Além das figuras de som e pausas, há outros símbolos que frequentemente se encontram numa pauta. Veremos em detalhe a ligadura de prolongação e o ponto de aumentação. \textbf{Ligadura}\textsc{[LOC]} de prolongação A ligadura de prolongação é uma linha curva que se emprega por cima ou por baixo das notas de música. A ligadura de prolongação aplica-se a duas ou mais notas iguais. Para notas diferentes usa-se a ligadura de expressão (ver abaixo). Quando se utiliza a ligadura de prolongação, só se toca a primeira nota, prolongando-se o som pelo valor da(s) nota(s) seguinte(s).

Por exemplo, duas mínimas unidas por uma ligadura de prolongação têm a mesma duração que uma semibreve; duas semínimas ligadas têm a mesma duração que uma mínima; etc. Ponto de aumentação Este ponto é colocado imediatamente à frente da figura que se quer aumentar. O ponto de aumentação adiciona à nota metade do seu valor. \textbf{Vejamos}\textsc{[MISC]} alguns exemplos:

Com a aplicação do ponto de aumentação, uma semibreve vale 6 tempos; uma mínima vale 3 tempos; e uma semínima vale um tempo e meio. Duplo ponto de aumentação Por vezes é necessária a aplicação de um segundo ponto. O duplo ponto de aumentação funciona como o primeiro ponto, isto é, aumenta metade do valor da figura que o precede. Todavia, neste caso, a figura precedente é o primeiro ponto de aumentação. Logo, o segundo ponto de aumentação aumenta metade do valor do primeiro ponto. \textbf{Vejamos}\textsc{[MISC]} o seguinte exemplo:

Uma mínima pontuada passa a valer três tempos e meio. Embora os exemplos apresentados tenham feito referência apenas a figuras de som, as pausas podem também ser pontuadas da mesma maneira.

Outras figuras Além das figuras musicais, da ligadura de prolongação e dos pontos de aumentação, existem muitos outros símbolos. Deixamos aqui alguns dos mais comuns: \textbf{Ligaduras}\textsc{[LOC]} A ligadura é um sinal de forma semicircular que se coloca acima ou abaixo das notas para ligar sons. Existem 3 tipos de ligadura: prolongação, expressão e de frase ou fraseado. A de prolongação, como já vimos, é a união de duas ou mais notas da mesma altura e mesmo nome. As durações das notas são somadas e ela é tocada como uma única nota. A ligadura de expressão liga duas notas de nomes diferentes. A ligadura de frase ou fraseado liga três ou mais notas de nomes diferentes. Legato Notas cobertas por este símbolo devem ser tocadas sem nenhuma interrupção como se fossem uma só.

Glissando Uma variação contínua de altura entre os dois extremos.

\textbf{Marca}\textsc{[PER]} de fraseado \textbf{Indica}\textsc{[MISC]} como as notas devem ser ligadas para formar uma frase. A execução varia de acordo com o instrumento.

\textbf{Tercina Condensa}\textsc{[PER]} três notas na duração que normalmente seria ocupada por apenas duas. Se as notas forem unidas por uma barra de ligação, as chaves ao lado do número podem ser omitidas. Grupos maiores podem ser formados e recebem o nome genérico de quiálteras ("que altera"), em que um certo número de notas é condensado numa duração inferior à soma das suas partes. Por exemplo, seis notas tocadas na duração que seria ocupada por quatro notas.

\textbf{Acorde Três}\textsc{[MISC]} ou mais notas tocadas simultaneamente. Se são tocadas apenas duas notas, dá-se o nome de intervalo.

\textbf{Arpejo}\textsc{[LOC]}, \textbf{Harpejo}\textsc{[LOC]} ou \textbf{Arpeggio Semelhante}\textsc{[LOC]} a um acorde, mas as notas não são tocadas simultaneamente, mas sim sucessivamente (uma de cada vez em sequência). \textbf{Marca}\textsc{[PER]} de respiração \textbf{Numa partitura vocal}\textsc{[MISC]}, indica o momento correto de fazer uma inspiração. Muitas vezes indica também um curtíssimo silêncio entre duas notas. Fermatta Uma nota sustentada indefinidamente, tendo a sua duração original prolongada ao gosto do executante. A fermata também pode aparecer sobre uma pausa, indicando uma suspensão, ou sobre a barra de compasso, indicando uma cesura (paragem total de som). \textbf{Compassos}\textsc{[MISC]} de espera \textbf{Marcação}\textsc{[LOC]} abreviada de pausa, indicando por quantos compassos deve-se manter a pausa.

\textbf{Ottava alta Ou}\textsc{[MISC]} oitava acima. Notas abaixo da linha pontilhada são tocadas uma oitava acima do escrito.

\textbf{Ottava Bassa}\textsc{[LOC]} Ou oitava abaixo. Notas abaixo da linha pontilhada são tocadas uma oitava abaixo do escrito.

\textbf{Quindicesima}\textsc{[LOC]} alta Notas abaixo da linha pontilhada são tocadas duas oitavas acima do escrito.

\textbf{QuindicesimaBassa Notas}\textsc{[MISC]} abaixo da linha pontilhada são tocadas duas oitavas abaixo do escrito.

No sistema musical ocidental utilizado existem 12 notas de música, constituindo a \textbf{Escala Cromática}\textsc{[LOC]}. Sete dessas notas são naturais e cinco são acidentes (ver \textbf{Acidentes}\textsc{[LOC]} a seguir). As 12 notas são as seguintes:

Cada nota encontra-se a meio-tom de distância da anterior e da seguinte. De referir que as únicas notas naturais que se encontram a meio-tom são Mi/Fá e \textbf{Si/Dó}\textsc{[MISC]}, constituindo os chamados meios-tons naturais. Os acidentes (as teclas pretas no piano) são conhecidos por dois nomes, relativos às notas naturais que as antecedem e sucedem. Assim sendo, \textbf{Dó\#}\textsc{[MISC]} é a nota que está meiotom acima de Dó e \textbf{Réb}\textsc{[MISC]} a nota que está meio-tom abaixo de \textbf{Ré --}\textsc{[MISC]} referindo-se, na prática, à mesma nota. Uma oitava (a extensão total da escala cromática) contém 12 meios tons, ou seja, 6 tons.

\textbf{Acidentes Deslocações}\textsc{[PER]} de tom ou acidentes são alterações efetuadas à nota de forma a mudar a sua frequência, tornando-a mais aguda ou mais grave. Existem 5 tipos de acidentes: o sustenido, o \textbf{bemol}\textsc{[MISC]}, o duplo sustenido, o duplo \textbf{bemol}\textsc{[MISC]} e o bequadro. São representados sempre antes do símbolo da nota cuja altura será modificada e depois do nome das notas, acordes e tonalidades. Um sustenido desloca a nota meio-tom acima (na escala), um duplo sustenido desloca o som um tom acima. Um \textbf{bemol}\textsc{[MISC]} desloca a nota meio-tom abaixo e o duplo \textbf{bemol}\textsc{[MISC]}

desloca o som um tom abaixo. O bequadro anula o efeito tanto dos bemóis como dos sustenidos. Por exemplo, pode-se dizer que um "Fá sustenido" (\textbf{Fá\#}\textsc{[MISC]}) é a mesma nota que um "\textbf{Sol bemol}\textsc{[MISC]}" (\textbf{Solb}\textsc{[LOC]}). Porém, devido às características de cada instrumento (e à própria disposição da escala), o timbre pode variar. Por exemplo, no caso da guitarra, um Dó tocado na segunda corda (Si), primeira posição, é equivalente a um Dó tocado na terceira corda (\textbf{Sol}\textsc{[PER]}) na quinta posição, embora o timbre seja diferente. Resumindo:

Uma vez que um sustenido ou \textbf{bemol}\textsc{[MISC]} tenha sido aplicado a uma nota, todas as notas de mesma altura manterão a alteração até ao fim do compasso. No compasso seguinte, todos os acidentes perdem o efeito e, se necessário, deverão ser aplicados novamente. Se desejarmos anular o efeito de um acidente aplicado imediatamente antes ou na armação de clave, devemos usar um bequadro, que faz a nota voltar à sua condição natural. No exemplo visto acima, podemos observar que a terceira nota do primeiro compasso também é sustenida, pois o acidente aplicado à nota anterior permanece válido e só é anulado pelo bequadro que faz a quarta nota voltar a ser um Lá natural. O segundo compasso é semelhante a não ser pelo acidente aplicado que é um \textbf{bemol}\textsc{[MISC]}. No terceiro compasso temos uma nota \textbf{Sol}\textsc{[PER]}, um \textbf{Lább}\textsc{[PER]} e um \textbf{Fá\#}\textsc{[MISC]}\#, o que faz com que as três notas

soem exatamente iguais, embora tenham nomes diferentes e ocupem posições também elas diferentes na clave. Alterações de precaução é o nome dado aos bequadros utilizados em situações que, em rigor, não seriam necessárias, tendo unicamente o fim de evitar dúvidas.

Na notação musical, um compasso é uma forma de dividir um trecho musical em pequenos fragmentos, com base em pulsações e pausas. Muitos estilos musicais tradicionais já presumem um determinado compasso: a valsa, por exemplo, centra-se no compasso 3/4 e o rock tipicamente usa os compassos 4/4 ou 12/8. Os compassos facilitam a execução musical, ao definir a unidade de tempo, a pulsação e o ritmo da composição -- ou de partes dela. Os compassos são formados na partitura a partir de linhas verticais desenhadas sobre a pauta (rever os símbolos apresentados na secção \textbf{Notação Musical}\textsc{[MISC]}). A soma dos valores temporais das notas e pausas dentro de um compasso deve ser igual à duração definida pela fórmula de compasso.

Fórmula de compasso Numa fórmula de compasso (ver figura abaixo), o denominador (número inferior) indica em quantas partes uma semibreve deve ser dividida para obtermos uma unidade de tempo. (Na notação atual, a semibreve é a medida com maior duração possível para ser atribuída a um tempo, sendo ela tomada como referência para as restantes durações.) O numerador (número superior) define quantas unidades de tempo o compasso contém.

Por exemplo, imaginemos o compasso "dois por quatro" do exemplo acima. A unidade de tempo tem duração de 1/4 da semibreve (uma semínima) e o compasso tem 2 unidades de tempo. Neste caso, uma mínima iria ocupar todo o compasso. Cada compasso pode ter qualquer combinação de notas e pausas, mas a soma de todas as durações nunca pode ser menor ou maior do que duas unidades de tempo (neste exemplo). A fórmula de compasso é escrita no inicio da composição ou de cada uma de suas secções (depois da armação de clave). Quando ocorre mudança de fórmula durante a música, essa mudança é escrita diretamente no compasso que tem a nova duração.

\textbf{Fórmula de compasso Indica}\textsc{[MISC]} qual o tipo de compasso em que a música está escrita.

\textbf{Tempo quaternário}\textsc{[MISC]} Este é o tempo mais usado e representa abreviadamente uma fórmula de 4/4.

\textbf{Tempo binário Indica}\textsc{[MISC]} um tempo de 2/2.

\textbf{Numerador}\textsc{[MISC]} Como já foi citado anteriormente, o número de cima (numerador) da fórmula de compasso indica a quantidade de unidades de tempo de cada compasso. Como este número indica a quantidade de pulsações em cada compasso, pode ser utilizado qualquer valor de numerador, desde que a estrutura do compasso esteja vinculada a uma ideia musical. Embora hajam alguns valores mais comuns, nada impede que um compositor utilize fórmulas com estruturas bastante complexas, principalmente em jazz, rock progressivo e música erudita contemporânea, onde fórmulas como 17/16, 19/16, 13/8 ou outros são comuns. Unidade de valor, de compasso e de tempo É importante não confundir unidade de valor, que é sempre a semibreve, com unidade de compasso ou de tempo. A unidade de compasso é o valor que preenche um compasso. Assim, a unidade de compasso de um 2/4 é a mínima, e num 2/2 é uma semibreve (isto é, o valor total do compasso). A unidade de tempo é o valor que preenche uma pulsação de qualquer compasso. A unidade de tempo num 2/4 é uma semínima e no 3/8 é uma colcheia.

\textbf{Classificações dos compassos}\textsc{[MISC]} Os compassos podem ser classificados de acordo com dois critérios: se levarmos em conta as notas que o compõem podemos dividi-los em simples e compostos. Se por outro lado considerarmos a métrica, eles podem ser binários, ternários, quaternários ou complexos.

Compasso simples é aquele em que cada unidade de tempo corresponde à duração determinada pelo denominador da fórmula de compasso. Por exemplo um compasso 2/4 possui duas pulsações com duração de 1/4 (uma semínima) cada. Os tipos mais comuns de compassos simples possuem 2 ou 4 no denominador (2/2, 2/4, 3/4 ou 4/4). Compasso binário \textbf{Célula}\textsc{[LOC]} rítmica formada por duas unidades de tempo: a pulsação é forte - fraca, ou seja, o primeiro tempo do compasso é forte e o segundo é fraco. Exemplos de binários simples são os compassos 2/8, 2/4, 2/2. Compasso ternário Métrica formada por três unidades de tempo. Exemplos de ternários simples são 3/4 ou 3/2. Compasso quaternário \textbf{Composto}\textsc{[PER]} por quatro unidades de tempo. Pode ser formada pela aglomeração de dois binários. A aglomeração pode ser notada quando o primeiro tempo é acentuado, segundo e quarto são fracos e o terceiro tem intensidade intermediária. São alguns exemplos de compasso quaternário simples 4/2, 4/4, 4/8, 4/16.

\textbf{Compasso composto}\textsc{[MISC]} \textbf{Compasso composto}\textsc{[MISC]} é aquele em que cada unidade de tempo é subdividida em três notas, cuja duração é definida pelo denominador da fórmula de compasso. Por exemplo, no compasso 6/8, o denominador indica que uma semibreve foi dividida em 8 partes (em colcheias) e o numerador indica quantas figuras preenchem o compasso, ou seja, o compasso é formado por 6 colcheias. No entanto, a métrica deste compasso é binária, ou seja, 2 pulsações por compasso. Por isso, a unidade de tempo não é uma colcheia, mas sim um grupo de três colcheias (ou uma semínima pontuada). Como cada pulsação é composta de três notas, este compasso é definido como composto. Obtém-se um compasso composto multiplicando um compasso simples pela fração de 3/2 -- ou seja, multiplicando o número de cima por 3 e o de baixo por 2. Por exemplo: o compasso 2/4 é um binário simples. Usando a fórmula (2/4)*(3/2) obtemos um 6/8 que corresponde a um binário composto. 3/4 é um ternário simples, e da mesma forma obtemos um 9/8 que corresponde a um ternário composto. 4/4 é um quaternário simples, que equivale a um 12/8, ou seja, um quaternário composto.

Compasso binário Alguns exemplos de binário composto são 6/4 6/8, 6/16, desde que haja divisão binária. Compasso ternário Alguns exemplos de compassos ternários compostos são 9/8, 9/16. Compasso quaternário São alguns exemplos de compasso quaternário composto 12/4, 12/8, 12/16.

Compasso complexo Existe uma característica auditiva que não nos permite ouvir compassos superiores a quatro tempos sem contá-los ou subdividi-los em valores menores. Por isso, os compassos acima de 4 tempos apresentam sempre uma subdivisão interna em partes menores ou iguais a 4 tempos. Quando alguns compositores utilizam compassos com métricas 5/4, 5/8, 7/8, 10/8, 11/8 e outras, trata-se sempre de aglomerações. No 5/4, por exemplo, trata-se da justaposição de um 2/4, seguido de um 3/4 (ou vice-versa). Outro exemplo é o 7/4 que pode ser formado por um 2/4, um 3/4 e outro 2/4, ou por um 4/4 e um 3/4.

\textbf{Cinética}\textsc{[LOC]} musical \textbf{Cinética Musical}\textsc{[MISC]} define a velocidade de execução de uma composição. Esta velocidade é chamada de andamento e indica a duração da unidade de tempo. O andamento é referido no início da música ou de um movimento e é indicada por expressões de velocidade em italiano, como \textbf{Allegro}\textsc{[ORG]} (rápido) ou \textbf{Addagio}\textsc{[LOC]} (lento). Junto ao andamento, pode ser indicada a expressão com que a peça deve ser interpretada, como com afecto, intensamente, melancólico, etc. \textbf{Vejamos}\textsc{[MISC]} abaixo uma lista de andamentos: \textbf{Gravíssimo:}\textsc{[MISC]} Menos de 40 batidas por minuto (bpm). Extremamente lento. Grave: de 40 a 48 bpm. Muito lento; grave; sério; demasiadamente vagaroso; \textbf{Largo}\textsc{[LOC]}: de 48 a 58 bpm. Lento, muito vagaroso; \textbf{Larghetto}\textsc{[PER]}: de 59 a 65 bpm. Um pouco mais rápido que o largo \textbf{Adágio}\textsc{[PER]}: de 66 a 72 bpm. \textbf{Devagar}\textsc{[PER]}; calmo; lentamente \textbf{Andante}\textsc{[LOC]}: de 73 a 80 bpm. Em passo tranquilo; andando \textbf{Andantino}\textsc{[PER]}: de 80 a 95 bpm. Um pouco mais rápido que o andante \textbf{Moderato}\textsc{[PER]}: de 96 a 104 bpm. Velocidade moderada; moderadamente. Allegretto: de 105 a 120 bpm. Mais rápido que o moderato e mais lento que allegro \textbf{Allegro}\textsc{[ORG]}: de 121 a 140 bpm. Depressa; rápido \textbf{Vivace}\textsc{[PER]}: de 141 a 168 bpm. \textbf{Vivo}\textsc{[ORG]}; com vivacidade; \textbf{Presto}\textsc{[MISC]}: de 169 a 180 bpm. Muito depressa; muito rápido \textbf{Prestíssimo}\textsc{[PER]}: de 181 a 208 bpm. O mais depressa possível Alguns exemplos de combinações de andamento com expressões: ? \textbf{Allegro moderato}\textsc{[MISC]} -- Moderadamente rápido. ? \textbf{Presto}\textsc{[MISC]} com fuoco -- Extremamente rápido e com expressão intensa. ? \textbf{Andante Cantabile -- Velocidade}\textsc{[MISC]} moderada e entoando as notas como que numa canção. ? \textbf{Adágio Melancolico}\textsc{[LOC]} -- Lento e melancólico

\textbf{Notações}\textsc{[LOC]} de variação de tempo: ? \textbf{Rallentando -- Indica}\textsc{[MISC]} que a execução deve tornar-se gradativamente mais lenta ? \textbf{Accelerando -- Indica}\textsc{[MISC]} que a execução deve tornar-se mais rápida. ? A tempo ou \textbf{Tempo primo -- Retorna}\textsc{[MISC]} ao andamento original. ? \textbf{Tempo rubato -- Indica}\textsc{[MISC]} que o músico pode executar com pequenas variações de andamento ao seu critério. \textbf{Marca}\textsc{[PER]} de metrónomo \textbf{Escrita}\textsc{[LOC]} no inicio da partitura, indica precisamente a duração de uma unidade de tempo (ou de uma pulsação) em batidas por minuto. Neste exemplo, a marca indica que 120 unidades de tempo (semínimas) ocupam um minuto (ou seja, que a pulsação é de 120 batidas por minuto - 120 BPM).

A intensidade das notas pode (e deve!) variar ao longo de uma música. A isso dá-se o nome de dinâmica. A intensidade é indicada em forma de siglas que indicam expressões em italiano sob a pauta. Certos símbolos e textos indicam ao intérprete a forma de executar a partitura, incluindo as variações de volume (dinâmica) e tempo (cinética), assim como a maneira correta de articular as notas e separá-las em frases (articulação e acentuação). \textbf{Vejamos}\textsc{[MISC]} os símbolos usados para anunciar a dinâmica pretendida: \textbf{Pianíssimo Execução}\textsc{[MISC]} muito suave. \textbf{Piano Suave}\textsc{[PER]}. \textbf{Mezzo}\textsc{[PER]} piano A intensidade é moderada, não tão fraca quanto o piano. \textbf{Mezzo}\textsc{[PER]} forte A intensidade é moderadamente forte. \textbf{Forte}\textsc{[LOC]} A intensidade é forte. \textbf{Fortíssimo}\textsc{[PER]} A intensidade é muito forte.

\textbf{Crescendo}\textsc{[PER]} Um crescimento gradual do volume. Esta marca pode ser estendida ao longo de muitas notas para indicar que o volume cresce gradualmente ao longo da frase musical. Diminuendo Uma diminuição gradual do volume. Pode ser estendida como o crescendo. Sforzando \textbf{Denota}\textsc{[LOC]} um aumento súbito de intensidade.

A figura abaixo mostra um solo de trompa da \textbf{Sinfonia}\textsc{[LOC]} número 5 de \textbf{Tchaikovsky}\textsc{[PER]}. Esta partitura apresenta várias marcas de dinâmica.

\textbf{Intervalo}\textsc{[PER]} é a relação entre as frequências de duas notas. São classificados quanto à simultaneidade ou não dos sons e à distância (altura) entre eles. Os intervalos são denominados pelo número de graus que contêm e qualificados pelo número de tons e meios-tons que possuem.

Tipos de intervalos \textbf{Intervalos Melódicos}\textsc{[ORG]} e \textbf{Harmónicos}\textsc{[LOC]} Na escala diatónica, a primeira classificação de um intervalo é quanto à ocorrência de simultaneidade na sua execução. Assim, o intervalo será melódico quando os sons aparecerem sucessivamente (um após o outro), ou harmónico, caso sejam executados simultaneamente (ao mesmo tempo). \textbf{Intervalos}\textsc{[MISC]} ascendentes ou descendentes Os intervalos podem também ser classificados como ascendentes, onde a segunda nota é mais aguda do que a primeira, ou descendentes, onde a segunda nota é mais grave do que a primeira. \textbf{Intervalos}\textsc{[MISC]} Simples e \textbf{Composto}\textsc{[PER]}s Os intervalos também podem ser simples ou compostos, dependendo da distância entre as notas. Um intervalo simples não excede uma oitava; um intervalo composto ultrapassa uma oitava. \textbf{Intervalo melódico Pode}\textsc{[MISC]} ser classificado quanto: ? À posição do segundo som em relação ao primeiro (ascendente ou descendente)

? À distância entre os dois sons. Será conjunto o intervalo que dista de um ou dois meios-tons (somente o intervalo de segunda) entre as notas e serão disjuntos todos os outros. \textbf{Intervalo harmónico}\textsc{[PER]} O intervalo harmónico pode ser classificado somente quanto à distância entre os dois sons. Uníssono Entre dois sons da mesma altura não há intervalo; a junção desses dois sons forma o uníssono.

\textbf{Nomes dos Intervalos}\textsc{[MISC]} Os nomes dos intervalos da escala diatónica são dados pela distância entre as notas, isto é, distância de segunda entre duas notas, terceira entre três, quarta, quinta, sexta, sétima, oitava, nona, etc.

De seguida, acrescenta-se o designativo que indica se a frequência entre os intervalos é mais ou menos consonante -- intervalo perfeito (justo), menor, maior, aumentado, diminuto, super-aumentado ou super-diminuto -- chamado também de "qualidade" do intervalo. Assim, temos os seguintes intervalos: \textbf{Tons}\textsc{[MISC]}:

Um intervalo menor, quando decrescido de um meio-tom, transforma-se num intervalo diminuto. Um intervalo maior, quando acrescido de um meio-tom, transforma-se num intervalo aumentado. Um intervalo diminuto, quando decrescido de um meio-tom, transforma-se num intervalo super diminuto. Um intervalo aumentado, quando acrescido de um meio-tom, transforma-se num intervalo super aumentado. Um intervalo perfeito, quando decrescido de um meio-tom, transforma-se num intervalo diminuto, e quando acrescido de um meio-tom, transforma-se num intervalo aumentado. Ou seja, por ordem ascendente (de meios tons), temos os seguintes intervalos: superdiminuto, diminuto, menor, maior, perfeito (utilizado em vez do menor e maior), aumentado e super aumentado. No caso de dois intervalos com a mesma distância, mas com nomes diferentes, como, por exemplo, a quarta aumentada e a quinta diminuta (ambos com 3 tons), ou a terceira diminuta e a segunda maior (ambos com 1 tom), dá-se o nome de intervalos inarmónicos. Na hora de identificar os intervalos, existe um meio mais racional e fácil de se saber a qualidade de um dado intervalo sem ter de contar o número de meios tons entre as notas. Basta estar atento ao facto de que, na escala diatónica ou natural, a distância entre todas as notas é de um tom, exceto entre as notas Mi e Fá e \textbf{Si e Dó}\textsc{[MISC]}, onde o intervalo é de um meio-tom - são os chamados meios-tons naturais. Uma vez localizados esses meios-tons naturais, basta levar em conta que: ? Nos intervalos de segunda e terceira, são MAIORES os que não possuem meio-tom, isto é, não "passam por" nenhum meio-tom natural. ? Nos intervalos de sexta e sétima, são MAIORES os que passam apenas por um meiotom natural. ? Os intervalos de quarta e quinta são todos PERFEITOS, com exceção do trítono (quarta aumentada, Fá para Si, ou quinta diminuta, \textbf{Si para Fá}\textsc{[MISC]} da oitava seguinte).

Inversão de intervalos Obtém-se a inversão de qualquer intervalo mudando a nota inferior para uma oitava acima ou mudando a nota superior para uma oitava abaixo. A inversão de intervalos funciona sempre num universo numérico de 9. Ora vejamos: O uníssono transforma-se em oitava; A segunda em sétima; A terceira em sexta; A quarta em quinta; A quinta em quarta; A sexta em terceira; A sétima em segunda; E a oitava em uníssono.

A sua qualificação, excetuando os intervalos justos (ou perfeitos), que conservam a primitiva, transforma-se no seu oposto: Os maiores ficam menores; Os menores ficam maiores; Os diminutos ficam aumentados; Os aumentados ficam diminutos. A inversão de intervalos pode ser particularmente útil no cálculo de intervalos entre duas notas distantes. Por exemplo, o intervalo entre \textbf{Dó e Sib}\textsc{[PER]} (ascendente) pode não ser imediatamente óbvio. Mas, utilizando a inversão (colocando o Dó uma oitava acima), ficamos com o intervalo \textbf{Sib}\textsc{[LOC]} e Dó. Sendo este uma \textbf{Segunda Maior}\textsc{[MISC]} (composta por um tom), a inversão mostra-nos facilmente que o intervalo original (\textbf{Dó e Sib}\textsc{[PER]}) é uma \textbf{Sétima Menor}\textsc{[LOC]} (composta por cinco tons).

Uma escala musical é uma sequencia ordenada de tons pela frequência vibratória de sons, (normalmente do som de frequência mais baixa para o de frequência mais alta), que consiste na manutenção de determinados intervalos entre as suas notas. Cada nota da escala denomina-se grau (primeiro grau, terceiro grau, etc.). Em \textbf{solfejo}\textsc{[LOC]}, as sílabas para representar as notas de quaisquer escalas são: \textbf{Dó, Ré}\textsc{[MISC]}, Mi, Fá, \textbf{Sol}\textsc{[PER]}, Lá, Si. Estas notas, no \textbf{Mundo Anglo-saxónico}\textsc{[MISC]}, são representadas pelas seguintes letras: C, D, E, F, G, A, B. A Lá

\textbf{Graus}\textsc{[ORG]} Em teoria musical, o grau determina a posição de uma nota musical em relação à primeira nota da escala diatónica. Cada grau e representado por um número romano e recebe um nome próprio, conforme o quadro seguinte: \textbf{Ordem 1º}\textsc{[ORG]} 2º 3º 4º 5º 6º 7º

No caso do grau VII, utilizam-se nomes específicos conforme o intervalo relativo criado com a tónica superior (grau VIII): chama-se subtónica quando é de um tom (escala menor natural), e sensível, quando o intervalo é de meio tom (escala maior, escala menor harmónica e escala menor melódica).

\textbf{Graus Modais}\textsc{[ORG]} Os graus modais, numa escala diatónica, são aqueles que pela sua relação com a tónica revelam o modo da escala -- maior ou menor. Estes são o 3º e 6º graus da escala. Se eles formarem respetivamente uma 3ª maior e uma 6ª maior, então a escala está no modo maior. Se, por outro lado, formarem uma 3ª menor e uma 6ª menor, então a escala está no modo menor (natural). \textbf{Graus}\textsc{[ORG]} conjuntos e disjuntos \textbf{Graus}\textsc{[ORG]} conjuntos são aqueles que se sucedem pela sua ordem numérica, como do 1º para o 2º. Os graus disjuntos são aqueles que não se sucedem, como do 4º para o 7º ou do 2º para o 6º. \textbf{Escala simples}\textsc{[MISC]} \textbf{Escala simples}\textsc{[MISC]} é a que não excede a 8ª justa/perfeita. As escalas começam e acabam com a tónica. Qualquer seguimento de notas, mesmo por graus conjuntos, que não comece e acabe com a tónica, é um fragmento de escala.

\textbf{Escala Cromática}\textsc{[LOC]} Como já vimos na página 13, a escala cromática contém 12 notas com intervalos de meios-tons entre elas. Esta escala foi criada no mundo ocidental através do estudo das frequências sonoras. A escala é formada pelas 7 notas padrão da escala diatónica acrescidas dos 5 tons intermédios (acidentes).

\textbf{Escala Diatónica do Modo Maior}\textsc{[MISC]} A 7 notas dispostas progressivamente (com a duplicação da primeira, perfazendo um total de oito notas) dá-se o nome de escala diatónica. As escalas podem começar por qualquer nota, tomando o nome da sua tónica (o primeiro grau). Uma escala que comece em Dó é a escala de Dó. Uma que comece em \textbf{Fá\#}\textsc{[MISC]} será a escala de \textbf{Fá\#}\textsc{[MISC]}.

Esta escala de dó maior é a escala modelo (pois não tem quaisquer acidentes). Uma escala que comece por qualquer outra nota tem de manter a mesma uniformidade ou estrutura desta escala modelo relativamente ao posicionamento de tons e meios-tons. Para se obter esta uniformidade, é indispensável recorrer-se à utilização de acidentes (sustenidos ou bemóis). Em todas as escalas diatónicas de modo maior devem encontrar-se dois meios-tons: um entre o 3º e 4º graus e outro entre o 7º e 8º graus. Todos os restantes graus encontramse a um tom de distância entre si. Seguem-se dois exemplos de escalas maiores que requerem a utilização de acidentes: \textbf{Sol Maior}\textsc{[PER]}

No primeiro exemplo (\textbf{Sol}\textsc{[PER]} maior), a alteração que tornou uniforme a sucessão de tons/meios-tons da escala foi empregue acidentalmente, ou seja, junto à nota alterada. Como já vimos, as alterações que se empregam na formação das escalas colocam-se na armação de clave, como se pode ver no segundo exemplo (escala de \textbf{Fá maior}\textsc{[MISC]}).

\textbf{Escala Diatónica do Modo Menor}\textsc{[MISC]} Quando se diz escala menor, entende-se escala diatónica do modo menor. Há três fórmulas na escala menor: natural, harmónica e melódica. \textbf{Fórmula Natural}\textsc{[MISC]} Na escala menor pela fórmula natural, os meios-tons encontram-se do 2º para o 3º e do 5º para o 6º graus. Entre os restantes graus, encontra-se um tom. No modo menor há uma escala que serve de modelo para todas as escalas menores. Essa escala é a de Lá menor, que é construída sem qualquer alteração constitutiva (acidentes).

\textbf{Escalas}\textsc{[LOC]} relativas: As escalas menores naturais têm sempre uma relativa maior, isto é, uma escala que se forma exatamente com o mesmo número de alterações constitutivas (o mesmo número de sustenidos ou bemóis na armação de clave). A relativa maior encontra-se uma 3ª menor acima da sua relativa menor. Por exemplo, \textbf{Dó maior}\textsc{[MISC]} é a relativa de Lá menor -- têm o mesmo número de alterações constitutivas (neste caso, nenhuma) e encontram-se a um tom e meio de distância (uma 3ª menor). Outros exemplos são: \textbf{Sol}\textsc{[PER]} maior como relativa de \textbf{Mi menor}\textsc{[MISC]} e \textbf{Fá\#}\textsc{[MISC]} menor como relativa de Lá maior. \textbf{Fórmula Harmónica}\textsc{[ORG]} Na escala menor pela fórmula harmónica, os meios-tons encontram-se do 2º para o 3º, do 5º para o 6º e do 7º para o 8º graus. Consequentemente, do 6º para o 7º grau, existe um tom e meio. Entre os restantes graus, encontra-se um tom. \textbf{Intervalos}\textsc{[MISC]} da escala harmónica: Além da diferença estabelecida pelos intervalos de 3ª e 6ª sobre a tónica, que na escala maior são maiores e na escala menor são menores, a contextura da escala menor é muito diferente da da escala maior. Na escala harmónica encontram-se os seguintes intervalos, que não existem na escala maior: - 2ª aumentada (que contém um tom e um meio tom): entre o 6º e o 7º grau; - 4ª diminuta (que contém um tom e dois meios tons): entre o 7º e o 3º grau; - 5ª aumentada (que contém quatro tons): entre o 3º e o 7º grau; - 7ª diminuta (que contém três tons e três meios-tons): entre o 7º e o 6º grau. \textbf{Fórmula Melódica}\textsc{[MISC]} Na escala menor pela fórmula melódica, os meios-tons encontram-se do 2º para o 3º e do 7º para o 8º graus na ordem ascendente, e do 6º para o 5º e do 3º para o 2º na ordem descendente (igual à fórmula natural). Entre os restantes graus, encontra-se um tom. A escala menor melódica é usada raramente.

\textbf{Modos}\textsc{[LOC]} gregos Na \textbf{Grécia antiga}\textsc{[LOC]}, as diversas formas de organizar os sons diferiam de região para região, consoante as tradições culturais e estéticas de cada uma delas. Assim, cada uma das regiões da antiga \textbf{Grécia}\textsc{[LOC]} deu origem a um modo (organização dos sons naturais) muito próprio, e que adaptou a denominação de cada região respetiva. Desta forma, aparecenos o modo dórico (da região da \textbf{Dória}\textsc{[LOC]}), o modo frígio (da \textbf{Frígia}\textsc{[LOC]}), etc. Os modos gregos são uma espécie de inversão da escala maior. Se tocarmos a escala de dó maior a partir da nota dó, teremos o modo dó jônio, que nada mais é do que a própria escala natural no seu estado fundamental (meios-tons entre 3º e 4º e 7º e 8º graus). Se tocarmos essa mesma escala a partir do segundo grau, a nota ré, teremos o modo ré dórico, e obteremos assim uma nova configuração da escala (meios-tons entre 2º e 3º e 6º e 7º graus) e consequentemente novos intervalos. Nos quadros abaixo vemos os 7 modos gregos, as suas estruturas e alguns exemplos: Jónio \textbf{Dórico}\textsc{[PER]} Frígio Lídio \textbf{Mixolídio}\textsc{[PER]} Eólio Lócrio

st -- meios-tons: 3º-4º; 7º-8º (esc. maior) T -- meios-tons: 2º- 3º; 6º- 7º T -- meios-tons: 1º- 2º; 5º- 6º st -- meios-tons: 4º- 5º; 7º- 8º T -- meios-tons: 3º-4º; 6º- 7º T -- meios-tons: 2º- 3º; 5º- 6º (esc. menor) T -- meios-tons: 1º- 2º; 4º- 5º.

\textbf{Armação}\textsc{[LOC]} de clave ou tonalidade não é mais do que a associação de sustenidos ou bemóis representados junto à clave, indicando a escala em que a música será expressa. Por exemplo, uma representação sem sustenidos ou bemóis será a escala de \textbf{Dó Maior}\textsc{[PER]} (ou a relativa Lá menor). Ao contrário dos acidentes aplicados ao longo da partitura, os sustenidos ou bemóis aplicados na clave duram por toda a peça ou até que uma nova armação seja definida (transposição, ou seja, mudança de tom).

Exemplo de uma armação com bemóis \textbf{Baixa}\textsc{[LOC]} a altura de todas as notas indicadas pelos bemóis nas posições indicadas junto a clave e as notas de mesmo nome em qualquer oitava. Os bemóis são acrescentados de acordo com a sequência inversa do círculo das quintas, ou seja, \textbf{Sib}\textsc{[LOC]}, \textbf{Mib}\textsc{[PER]}, \textbf{Láb}\textsc{[PER]}, \textbf{Réb}\textsc{[MISC]}, \textbf{Solb}\textsc{[LOC]}, \textbf{Dób}\textsc{[MISC]} e \textbf{Fáb}\textsc{[ORG]}. Exemplo de uma armação com sustenidos \textbf{Eleva}\textsc{[LOC]} a altura de todas as notas indicadas pelos sustenidos nas posições indicadas junto a clave e as notas de mesmo nome em qualquer oitava. Os sustenidos são acrescentados de acordo com a sequência do círculo das quintas, ou seja \textbf{Fá\#}\textsc{[MISC]}, \textbf{Dó\#}\textsc{[MISC]}, \textbf{Sol\#}\textsc{[MISC]}, \textbf{Ré\#}\textsc{[MISC]}, Lá\#, \textbf{Mi\#}\textsc{[MISC]} e \textbf{Si\#}\textsc{[MISC]}.

Olhando para a armação de clave, imediatamente sabemos em que tonalidade se encontra a peça: ? Se a armação não tiver acidentes, encontra-se em \textbf{Dó maior}\textsc{[MISC]} (ou Lá menor). ? Se tiver um \textbf{bemol}\textsc{[MISC]}, a tonalidade é \textbf{Fá maior}\textsc{[MISC]} (ou \textbf{Ré menor}\textsc{[MISC]}). ? Com vários bemóis, o penúltimo \textbf{bemol}\textsc{[MISC]} indica a tonalidade. ? Se tiver sustenidos, ao último sustenido sobe-se meio-tom para descobrir a tonalidade.

\textbf{Método}\textsc{[PER]} de cálculo das tonalidades \textbf{Teoricamente}\textsc{[LOC]}, existem 34 tonalidades diferentes (17 para escalas maiores e 17 para escalas menores). A maioria dos alunos de teoria deverá memorizar todas. Felizmente, usando o método de cálculo das tonalidades, é apenas necessário memorizar 7. Neste método, a cada tonalidade é atribuído um valor numérico baseado no número e tipo de acidentes. Os sustenidos são positivos, os bemóis negativos. \textbf{Dó maior}\textsc{[MISC]} não tem acidentes, logo o seu valor numérico é 0. \textbf{Ré maior}\textsc{[MISC]} tem 2 sustenidos, logo o seu valor é 2. \textbf{Mi maior}\textsc{[MISC]} tem 4 sustenidos, tendo um valor de 4. \textbf{Fá maior}\textsc{[MISC]} tem um \textbf{bemol}\textsc{[MISC]}, logo equivale a -1 (os bemóis assumem um valor negativo). \textbf{Sol}\textsc{[PER]} maior tem um sustenido, obtendo um valor numérico de 1. Lá maior tem 3 sustenidos, cujo valor numérico é de 3. Finalmente, \textbf{Si maior}\textsc{[MISC]} tem 5 sustenidos, recebendo um valor numérico de 5. Estas 7 escalas e respetivos valores deverão ser memorizados para que se possa utilizar este método com as restantes tonalidades. Comparando agora \textbf{Dób maior}\textsc{[MISC]}, \textbf{Dó maior}\textsc{[MISC]} e \textbf{Dó\# maior}\textsc{[MISC]}, vamos perceber como calcular os acidentes relativamente a todas as escalas maiores. Assim, se começarmos com \textbf{Dó maior}\textsc{[MISC]} e subtrairmos 7, acabamos em \textbf{Dób maior}\textsc{[MISC]} (que tem 7 bemóis). Se começarmos em \textbf{Dó maior}\textsc{[MISC]} e adicionarmos 7, chegamos a \textbf{Dó\# maior}\textsc{[MISC]}. Estes dois exemplos mostram a relação numérica entre as escalas: uma "escala \textbf{bemol}\textsc{[MISC]}" (7) está sempre 7 valores abaixo da "escala natural" (0) que está 7 valores abaixo da "escala sustenido" (+7).

Exemplos: Se quisermos encontrar \textbf{Mib}\textsc{[PER]} maior, começamos com \textbf{Mi maior}\textsc{[MISC]}, que tem um valor numérico de 4 (4 sustenidos). Para converter para \textbf{Mib}\textsc{[PER]} maior, subtraímos 7 àquele valor. O resultado é -3 (4 - 7 = -3). Assim sendo, \textbf{Mib}\textsc{[PER]} maior tem 3 bemóis. Como outro exemplo, vejamos \textbf{Fá\#}\textsc{[MISC]} maior. \textbf{Fá maior}\textsc{[MISC]} tem um \textbf{bemol}\textsc{[MISC]} (-1). Para converter para \textbf{Fá\#}\textsc{[MISC]} maior, adicionamos 7, sendo o resultado 6. Logo, \textbf{Fá\#}\textsc{[MISC]} maior tem 6 sustenidos. Para compreender agora as escalas menores (na fórmula natural), vamos comparar \textbf{Dó maior}\textsc{[MISC]} (que não tendo acidentes tem um valor de 0) com \textbf{Dó menor}\textsc{[MISC]} (com 3 bemóis, tendo um valor de -3). Assim, para converter uma escala maior na sua paralela menor, basta subtrair 3. \textbf{Dó Maior}\textsc{[PER]}: 0 Dó \textbf{Menor}\textsc{[PER]}: -3 Exemplos: \textbf{Vamos calcular}\textsc{[MISC]} \textbf{Ré menor}\textsc{[MISC]}. Começamos com \textbf{Ré maior}\textsc{[MISC]}, que tem um valor de 2. Em seguida, subtraímos 3. O resultado é -1. Assim, \textbf{Ré menor}\textsc{[MISC]} tem um \textbf{bemol}\textsc{[MISC]}. Agora vejamos \textbf{Fá menor}\textsc{[MISC]}. \textbf{Fá maior}\textsc{[MISC]} tem um valor de -1 (um \textbf{bemol}\textsc{[MISC]}). Subtraímos 3 e chegamos a um valor de -4. Logo, \textbf{Fá menor}\textsc{[MISC]} tem 4 bemóis. Algumas tonalidades exigem duas conversões. Por exemplo, vamos calcular \textbf{Dó\# menor}\textsc{[MISC]}. Começamos com \textbf{Dó maior}\textsc{[MISC]}, que tem um valor de 0. Depois adicionamos 7 para chegar a \textbf{Dó\# maior}\textsc{[MISC]}. Finalmente, subtraímos 3 para encontrar \textbf{Dó\# menor}\textsc{[MISC]}. O resultado é 4 (0+73=4), o que significa que \textbf{Dó\# menor}\textsc{[MISC]} tem 4 sustenidos.

Em música, acorde é a escrita ou execução de três ou mais notas simultaneamente -- para alguns teóricos, o acorde só se forma a partir de três ou mais notas, reservando a palavra intervalo para a execução de duas notas. Os acordes são formados a partir da nota mais grave, onde são acrescentadas as outras notas constituintes.

\textbf{Tríades Tríade}\textsc{[LOC]} é um acorde de 3 notas criado a partir de uma escala, normalmente a diatónica, com a sobreposição de duas terceiras (maiores ou menores). As suas três notas constituintes são a fundamental (nota mais grave e que dá o nome ao acorde), a 3ª (também chamada nota modal, que determina o carácter do acorde -- maior ou menor) e a 5ª. Existem 4 tipos de tríades possíveis de serem montadas a partir das escalas diatónicas maior e menor. São elas: diminuta, menor, maior e aumentada. Tríade: Diminuta

Quatríades Quatríades são acordes constituídos por 4 notas. Vamos abordar vários tipos de quatríades, entre elas os \textbf{Acordes Suspensos}\textsc{[LOC]}, \textbf{Acordes de Sexta}\textsc{[LOC]}, \textbf{Acordes de Sétima}\textsc{[LOC]} e \textbf{Extensões}\textsc{[LOC]}.

\textbf{Acordes Suspensos}\textsc{[LOC]} Um acorde suspenso, ou \textbf{sus chord}\textsc{[ORG]}, é um acorde na qual a terceira é substituída pela segunda ou quarta. Isto produz dois acordes principais: o acorde de segunda suspensa (sus2) e o de quarta suspensa (sus4). Os acordes C sus2 e \textbf{Csus4}\textsc{[MISC]}, por exemplo, consistem nas notas C D G e C F G (ver figura abaixo) respetivamente. Existe ainda um terceiro tipo de acorde suspenso, no qual tanto a segunda como a quarta estão presentes, resultando, por exemplo, num acorde de Dó com as notas C D F G.

O nome suspenso deriva de uma técnica na qual uma progressão melódica, antes de chegar a uma nota harmonicamente estável, ficava momentaneamente suspensa na nota anterior. Tal procedimento resultava numa dissonância inesperada que poderia posteriormente ser resolvida pelo surgimento da nota deslocada. Na teoria da musica tradicional, a inclusão de uma terceira juntamente com a segunda ou quarta negaria a suspensão e, dessa forma, tais acordes seriam chamados de nona (\textbf{added ninth}\textsc{[MISC]}) ou décima primeira (\textbf{added eleventh}\textsc{[MISC]}). No exemplo abaixo, podemos ver que no primeiro caso temos um acorde suspenso, pois não tem a terceira (E), enquanto que no segundo caso, devido à presença da terceira, temos um acorde de tom adicionado (add9).

\textbf{Acordes de Sexta}\textsc{[LOC]} A acorde de sexta maior (ex: C6) é de longe o tipo mais comum de acorde de sexta. É composto por uma tríade maior acrescida de uma sexta maior a partir da tónica. Por exemplo, o acorde C6 contém as notas C-E-G-A. Utilizando uma tríade menor, temos, por exemplo, o acorde \textbf{Cm6}\textsc{[LOC]}, composto pelas notas C-Eb-G-A. Na notação musical, em cada um destes acordes assume-se que a sexta é maior em relação à tónica. Porém, um intervalo de sexta menor pode ser indicado na notação do acorde da seguinte forma: Cm(m6) ou \textbf{Cm m6}\textsc{[MISC]}.

\textbf{Acordes de Sétima}\textsc{[LOC]} Um acorde de sétima consiste numa tríade, à qual é adicionado um intervalo de sétima acima da tónica. Uma variedade de sétimas pode ser adicionada a uma variedade de tríades, resultando em diferentes tipos de acordes de sétima, como descrito abaixo. Porque todas as sétimas são intervalos dissonantes, qualquer acorde de sétima é dissonante. Por outras palavras, um acorde de sétima é mais tenso do que um acorde maior ou menor. Todavia, alguns estilos musicais, como o jazz, utilizam abundantemente este tipo de acordes e trata-os como consonantes. \textbf{Acorde de Sétima Dominante}\textsc{[MISC]} Quando um acorde é denominado simplesmente como um acorde de sétima, usualmente significa que é um acorde de sétima dominante. Este acorde é construído tendo por base uma tríade maior e uma sétima menor (por exemplo, C-E-G-Bb, formando o acorde de C7). É chamado de sétima dominante porque o \textbf{acorde V}\textsc{[MISC]} (quinto grau da escala), ou acorde dominante, é o único acorde maior na escala maior que naturalmente contém uma sétima menor -- ver figura abaixo. É considerado como o mais importante de todos os acordes de sétima.

G7 (\textbf{Sol de Sétima Dominante}\textsc{[MISC]}) \textbf{Acordes de Sétima Maior e}\textsc{[PER]} \textbf{Menor}\textsc{[PER]} Enquanto o acorde de sétima dominante é tipicamente construído no quinto grau (dominante) de uma escala maior, os acordes maiores de sétima são habitualmente construídos no primeiro ou quarto grau da escala. É composto por uma tríade maior e uma sétima maior. Os acordes menores de sétima surgem no segundo, terceiro ou sexto graus de uma escala maior. É constituído por uma tríade menor e uma sétima menor.

Acorde de \textbf{Sétima Semi-Diminuta}\textsc{[MISC]} Um acorde de sétima semi-diminuta é considerado "semi-diminuto" porque é construído tendo por base uma tríade diminuta, à qual se acrescenta uma sétima menor, em vez de uma sétima diminuta. Como exemplo na tonalidade de Dó, temos as notas B D F e A.

Si de \textbf{Sétima Semi-Diminuta}\textsc{[MISC]} -- \textbf{Bø7}\textsc{[MISC]}. \textbf{Acorde de Sétima Diminuta}\textsc{[MISC]} Um acorde de sétima diminuta constrói-se sobrepondo 3 intervalos de terceira menor (ex.: B \textbf{D F Ab}\textsc{[MISC]}) Resumindo, os acordes de sétima mais comuns são: \textbf{Sétima Dominante}\textsc{[MISC]} -- 1 3 5 7b \textbf{Sétima Maior}\textsc{[MISC]} -- 1 3 5 7 \textbf{Sétima Menor}\textsc{[LOC]} 1 3b 5 7b \textbf{Sétima Semi-Diminuta}\textsc{[MISC]} -- 1 3b 5b 7b \textbf{Sétima Diminuta}\textsc{[MISC]} -- 1 3b 5b 7bb

\textbf{Extensões}\textsc{[LOC]} (extended chords) são tríades com terceiras acrescentadas acima da sétima, formando a nona, décima primeira e décima terceira. Depois da décima terceira, qualquer terceira que se acrescente iria duplicar notas já presentes no acorde. Ou seja, com estas extensões, todas as notas da tonalidade estão já presentes, de modo que acrescentar novos intervalos não iria acrescentar novos tons. Estes acordes só podem ser construídos usando notas situadas fora das sete presentes na escala diatónica, isto é, para além de uma oitava. Abaixo apresenta-se um quadro resumo dos principais acordes.

\textbf{Notação}\textsc{[LOC]} Os acordes aparecem notados com a primeira letra maiúscula quando se trata de um acorde maior: Dó - dó maior Ré - ré maior \textbf{Etc}\textsc{[LOC]}. Os acordes aparecem notados com a primeira letra minúscula quando se tratar de um acorde menor: sol - sol menor lá - lá menor etc. Quando forem acordes diminutos ou aumentados, escreve-se dim ou aum após o nome do acorde: dó dim - dó diminuto ré aum - ré aumentado

\textbf{Inversão de acordes}\textsc{[MISC]} Além das formas indicadas acima, é possível indicar acordes em que a sequência das notas é invertida e uma das notas mais agudas é usada como baixo. Para indicar a inversão, utiliza-se a mesma notação acima, indicando qual das notas deve ser o baixo do acorde, separada por uma barra. A primeira inversão é obtida movendo a nota fundamental oitava acima, passando a segunda nota a ser o baixo. A segunda inversão é realizada, movendo o baixo da primeira inversão uma oitava para cima. Por exemplo, se o acorde de \textbf{Sol}\textsc{[PER]} maior (G, B, D) for invertido uma vez, teremos G/B (B, D, G da oitava acima). Na segunda inversão o acorde é G/D (D, G da oitava acima, B da oitava acima). A terceira inversão não é indicada pois é igual ao acorde original, a única diferença sendo que todas as notas são tocadas uma oitava acima.

CIFRAS E TABLATURAS \textbf{Cifra}\textsc{[LOC]} é um sistema de notação musical usado para indicar através de símbolos gráficos ou letras os acordes a serem executados por um instrumento musical (como por exemplo uma guitarra). São utilizadas principalmente na música popular, acima das letras ou partituras de uma composição musical, indicando o acorde que deve ser tocado em conjunto com a melodia principal ou para acompanhar o canto. As principais cifras são grafadas: A: nota lá ou acorde de \textbf{Lá Maior B}\textsc{[MISC]}: nota si ou acorde de \textbf{Si Maior C}\textsc{[MISC]}: nota dó ou acorde de \textbf{Dó Maior D}\textsc{[PER]}: nota ré ou acorde de \textbf{Ré Maior}\textsc{[PER]} E: nota mi ou acorde de \textbf{Mi Maior F}\textsc{[MISC]}: nota fá ou acorde de \textbf{Fá Maior G}\textsc{[PER]}: nota sol ou acorde de \textbf{Sol Maior}\textsc{[PER]} Os acordes menores são grafados pelas letras acima, acompanhados da letra "m" minúscula. Ex: Cm indica um acorde de \textbf{Dó menor}\textsc{[MISC]}. Há outras alterações quando se utilizam quatríades ou intervalos dissonantes. Ex: \textbf{Cm7}\textsc{[MISC]} indica acorde de \textbf{Dó menor}\textsc{[MISC]} com sétima. A tablatura é uma notação que representa como colocar os dedos num instrumento (nos trastes de uma guitarra, por exemplo) em vez das notas, permitindo aos músicos tocar o instrumento sem formação especializada. Esta notação tornou-se comum na partilha de músicas pela \textbf{Internet}\textsc{[MISC]}.

\textbf{CONCLUSÃO Relembrando}\textsc{[MISC]}? ? A música escreve-se em \textbf{Pautas}\textsc{[LOC]} ou \textbf{Partituras}\textsc{[LOC]}, compostas por 5 linhas e 4 espaços. Nestas partituras surgem os mais variados símbolos: claves, fórmulas de compasso, figuras musicais, sinais de dinâmica e interpretação, etc. Ao conjunto destes símbolos dá-se o nome de \textbf{Notação Musical}\textsc{[MISC]}; ? As \textbf{Clave}\textsc{[LOC]}s são símbolos colocados no início da partitura e que, através da sua posição na mesma, indicam quais as notas de cada linha; ? Existem 7 \textbf{Figuras Musicais}\textsc{[MISC]} utilizadas atualmente (da semibreve à semifusa), todas com durações diferentes. Vimos também as suas respetivas pausas, tal como outros símbolos -- \textbf{Pontos de Aumentação}\textsc{[MISC]}, \textbf{Ligaduras de Prolongação}\textsc{[LOC]} e outros; ? São 12 as \textbf{Notas Musicais}\textsc{[MISC]} que compõem o sistema que utilizamos. Os \textbf{Acidentes}\textsc{[LOC]} (sustenidos e bemóis) são símbolos que, por vezes, são aplicados às 7 notas naturais, revelando as 12 notas existentes; \textbf{? Um Compasso}\textsc{[MISC]} é uma forma de dividir a música, podendo ser composto por 2 unidades de tempo (binário), 3 unidades de tempo (ternário) ou 4 unidades de tempo (quaternário). Existem compassos simples, compostos e complexos; ? Na música existem vários tipos de \textbf{Andamento}\textsc{[LOC]}, dos mais lentos aos mais rápidos, tal como outros termos que indicam alterações de velocidade na interpretação da peça; ? Existe também um conjunto de símbolos que nos indicam as alterações a levar a cabo na \textbf{Expressão e Intensidade}\textsc{[ORG]} da música. ? Os \textbf{Intervalos}\textsc{[MISC]} são a forma de medir a distância entre as notas, assumindo várias denominações e qualificações. A inversão de intervalos facilita a sua identificação; \textbf{? Nas Escalas Musicais}\textsc{[MISC]}, percebemos que existem a escala cromática, a escala diatónica do modo maior, a escala diatónica do modo menor (natural, harmónica e melódica) e os modos gregos (\textbf{Dórico}\textsc{[PER]}, \textbf{Mixolídio}\textsc{[PER]} e outros); \textbf{? A Tonalidade}\textsc{[MISC]} refere-se à armação de clave, onde são colocados os acidentes constitutivos da música; ? Há vários tipos de \textbf{Acordes}\textsc{[LOC]}, como por exemplo tríades e quatríades, podendo ser apresentados de forma invertida; \textbf{? Cifras}\textsc{[PER]} e \textbf{Tablaturas}\textsc{[LOC]} são formas alternativas de escrever o tema.


\begin{thebibliography}{1}
\bibitem{manual_lxpro}
LXPRO.
\textit{Manual de Teoria Musical}.
Documento PDF. Acedido em: março de 2026.
\end{thebibliography}

\end{document}
